%% kapa: exam=1 round=22 year=2011 period=1 form=A subject=economics subject_ko=경제원론 seq=39 number=79
%% source: Q-Net 감정평가사 기출문제 (cst00309) artlSeq=5036790 file=extracted/1008/5036790/2011년 제22회 감정평가사 1차시험 기출문제/1교시.hwp
%% conversion: kordoc 4.24.51 -> tools/gen1.py (auto); review=proofread-sonnet; flags=none
\begin{question}{79}
독점기업이 시장을 A, B로 구분하여 가격차별을 통해 이윤을 극대화하고 있다. 독점기업의 한계비용은 생산량과 관계없이 10으로 일정하고 현재 A, B 두 시장의 수요의 가격탄력성은 각각 2와 3이다. A, B 두 시장에서 독점기업이 설정하는 가격은?
\begin{choices}
  \item A: 30, B: 20
  \item A: 20, B: 15
  \item A: 15, B: 10
  \item A: 20, B: 30
  \item A: 25, B: 30
\end{choices}
\end{question}
\answer{2}
